\documentclass[11pt]{article}

\usepackage[a4paper,margin=1in]{geometry}
\usepackage{amsmath,amssymb,mathtools}
\usepackage[hidelinks]{hyperref}

\title{Ideas Imported from the STFT Nicola--Tilli Stability Paper}
\author{}
\date{}

\begin{document}
\maketitle

This note records what can be imported from
\texttt{s00222-024-01248-2-3.pdf}, the 2024 Inventiones paper

\begin{quote}
Gómez--Guerra--Ramos--Tilli, \emph{Stability of the Faber--Krahn inequality for
the short-time Fourier transform}.
\end{quote}

The paper proves a sharp quantitative stability version of the Nicola--Tilli
Faber--Krahn theorem for the short-time Fourier transform. It is not directly a
torsion result, but several mechanisms are relevant to \texttt{plan2.md}.

\section{The Structural Dictionary}

In the STFT/Fock setting, the central object is
\[
u_F(z)=|F(z)|^2e^{-\pi |z|^2}
\]
and its superlevel sets
\[
A_t=\{u_F>t\}.
\]
The rearranged profile \(u^*(s)\) satisfies a one-dimensional convexity
inequality. The optimizer profile is \(e^{-s}\). Stability is obtained by
quantifying the area between \(u^*(s)\) and \(e^{-s}\), and then pushing the
information back to sets.

For torsion, the analogous objects are
\[
u=u_\Omega,\qquad E_t=\{u>t\},\qquad v(s)=u_\Omega^*(s),
\]
and the optimizer profile is the ball profile
\[
b(s)=u_B^*(s)
=
\frac{R^2-(s/\omega_n)^{2/n}}{2n},
\qquad |B|=\omega_n R^n.
\]
The torsion level-set inequality gives
\[
v'(s)\ge b'(s).
\]
Thus
\[
h(s):=b(s)-v(s)
\]
is nonnegative and nonincreasing, and
\[
\int_0^{|\Omega|}h(s)\,ds
=
\int_Bu_B-\int_\Omega u_\Omega
=
2(E(\Omega)-E(B)).
\]
This is the exact torsion analogue of the STFT paper's ``area between profiles''
method.

\section{Immediate Lemma: Profile Deficit Controls Level Heights}

Since \(h\) is nonincreasing,
\[
0\le h(s)\le \frac{1}{s}\int_0^s h(r)\,dr
\le
\frac{2(E(\Omega)-E(B))}{s}.
\]
Therefore, for \(s\in[\theta|\Omega|,|\Omega|]\),
\[
|u_B^*(s)-u_\Omega^*(s)|
\le
\frac{2(E(\Omega)-E(B))}{\theta|\Omega|}.
\]
This gives quantitative control of near-boundary level heights. If \(s=L-\eta\)
and \(E(\Omega)-E(B)\ll \eta\), then the level
\[
t_\eta:=u_\Omega^*(L-\eta)
\]
is positive and comparable to the corresponding ball level. This selects a
boundary-layer set
\[
E_{t_\eta}=\{u>t_\eta\}
\]
with exactly
\[
|\Omega\setminus E_{t_\eta}|=\eta.
\]
This should be combined with the weighted level-set deficit identity to find
near-boundary levels that are also geometrically good.

\section{Superlevel Replacement}

The STFT proof does not try to control an arbitrary set directly. It first
replaces the set \(\Lambda\) by the matched superlevel set
\[
A_\Lambda=\{u_F>u_F^*(|\Lambda|)\}.
\]
Then it proves:
\begin{enumerate}
\item the deficit controls the function's distance to the optimizer;
\item this makes \(A_\Lambda\) regular and close to a ball;
\item a transport argument controls \(|\Lambda\Delta A_\Lambda|\);
\item hence \(\Lambda\) itself is close to a ball.
\end{enumerate}

For torsion, the analogous move is:
\begin{enumerate}
\item choose \(E_t=\{u_\Omega>t\}\) with prescribed volume \(s\);
\item prove \(E_t\) is close to a ball using the level-set identity;
\item use
\[
\mathcal A(\Omega)
\le
C\frac{|\Omega\setminus E_t|}{|\Omega|}
+C\mathcal A(E_t)
\]
to transfer the result back to \(\Omega\).
\end{enumerate}
The point of the profile-gap lemma is that it chooses near-boundary levels with
controlled height, while the exact level-set identity supplies averaged control
of \(D_H(t)+D_I(t)\).

\section{Transport Idea for the Set Comparison}

In the STFT paper, set stability is proved by transporting mass from
\(A_\Lambda\setminus\Lambda\) to \(\Lambda\setminus A_\Lambda\). Points moved
farther away from the center lose a definite amount of Gaussian weight, so the
concentration deficit controls the amount of badly moved mass.

For torsion, a literal radial transport is unavailable for arbitrary domains.
But for selected minimizers or near-ball sets, a normal-transport analogue should
be possible:
\begin{itemize}
\item compare a near-boundary level \(E_t\) with \(\Omega\) along the flow lines
of \(\nabla u\), or along the free-boundary normal when \(\Omega\) is already a
selected minimizer;
\item use the nondegeneracy
\[
c\operatorname{dist}(x,\partial\Omega)
\le u(x)\le
C\operatorname{dist}(x,\partial\Omega)
\]
near the free boundary to convert height \(t\) into boundary-layer thickness;
\item use density/perimeter bounds to estimate the symmetric difference.
\end{itemize}

This is one reason the hybrid route looks more promising than the pure
level-set route: the selected minimizers have exactly the free-boundary
structure needed to make this transport quantitative.

\section{Geometry of Superlevel Sets}

The STFT proof obtains star-shapedness, smoothness, and convexity of relevant
superlevel sets from closeness of \(F\) to the Gaussian and analyticity.

For torsion, the replacement should be:
\begin{itemize}
\item for arbitrary domains, no comparable regularity is available;
\item for selected minimizers, Alt--Caffarelli regularity gives \(C^{1,\gamma}\)
free boundary once flatness is known;
\item after entering the graph regime, Schauder estimates give \(u_\Omega\)
close to \(u_B\), and then regular near-boundary level sets follow from the
implicit function theorem.
\end{itemize}
So the STFT ``regular superlevel geometry'' mechanism should be imported only
after the BDV selection step.

\section{Variational/Fuglede-Style Route}

The STFT paper also proves stability by a variational second-variation argument
for
\[
K[F]=\frac{I_F(s)}{\|F\|^2},
\]
where \(I_F(s)\) is the integral of \(u_F\) over its superlevel set of measure
\(s\). Around the optimizer \(F\equiv1\), after removing neutral modes, the
second variation is uniformly negative.

The torsion analogue would be to study, near the ball,
\[
\mathcal K_s(\Omega)
:=
\int_{\{u_\Omega>u_\Omega^*(s)\}}u_\Omega\,dx
\]
under nearly spherical perturbations of \(\Omega\). The expected result is a
second-variation estimate of the form
\[
\mathcal K_s(B)-\mathcal K_s(\Omega)
\gtrsim_s
\|\varphi\|_{H^{1/2}(\partial B)}^2
\]
after imposing volume and barycenter orthogonality. This would be a
level-set-sensitive analogue of BDV's nearly spherical expansion for \(E\).

This is probably harder than the profile-gap lemma, but it could produce a
direct near-ball stability theorem for the level-set route.

\section{Revised Route for \texttt{plan2.md}}

The most promising revised route is:
\begin{enumerate}
\item Prove the finite-perimeter level-set identity.
\item Add the profile-gap lemma \(0\le b(s)-v(s)\le 2\delta_E/s\).
\item Choose boundary-layer volumes \(s=L-\eta\), with
\(\delta_E\ll\eta\ll1\).
\item Use the weighted deficit identity to find one such level with small
\(D_H+D_I\).
\item Apply quantitative isoperimetry to \(D_I\), and the weighted variance
identity to \(D_H\).
\item For arbitrary domains this still leaves a regularity gap; for selected
minimizers, use the BDV free-boundary estimates to close it.
\item Transfer from \(E_t\) back to \(\Omega\) by the boundary-layer asymmetry
lemma.
\end{enumerate}
This is a genuine import from the STFT stability paper: first stabilize the
matched superlevel set, then transfer stability back to the original set.

\section{Analytic/Native-Space Refinement}

The GGRT proof benefits from living in the Bargmann--Fock space. For torsion,
there is no global entire structure for arbitrary domains, but there is a
native harmonic substitute:
\[
h_\Omega(x):=u_\Omega(x)+\frac{|x|^2}{2N},
\qquad
\Delta h_\Omega=0.
\]
For a ball, \(h_B\) is constant. For a nearly spherical domain, after pulling
back to the ball, the first variation of \(h_\Omega-h_B\) is the harmonic
extension of the boundary graph. This is the natural space in which to try a
GGRT-style second variation.

In dimension two there is an even more literal analytic object:
\[
F_\Omega(z):=\partial_z u_\Omega(z)+\frac{\bar z}{4}.
\]
Since \(\partial_{\bar z}\partial_z u=-1/4\), the function \(F_\Omega\) is
holomorphic in \(\Omega\), and \(F_B\equiv0\) for disks centered at the origin.
For selected smooth near-disks, pulling \(F_\Omega\) back to the unit disk could
put the problem into a Bergman/Hardy-type setting, closer in spirit to the
Fock-space proof.

The concrete calculation to try is the second variation at the ball of
\[
\mathcal K_s(\Omega)
=
\int_{\{u_\Omega>u_\Omega^*(s)\}}u_\Omega\,dx.
\]
If the resulting quadratic form is diagonal in spherical harmonics and has a
spectral gap after removing the \(k=0\) and \(k=1\) modes, this would give a
native torsion analogue of the GGRT variational proof.

\end{document}
